% Figure : cycle autocohérent DFPT pour les phonons (fig:scf_phonon)
% Même style que le diagramme SCF électronique.
% Notation : \Delta \equiv \Delta_{\qv\alpha\mu} dans les boîtes pour alléger.
\begin{tikzpicture}[
    box/.style={rectangle, draw=black, rounded corners,
                fill=gray!15, minimum width=6.5cm,
                minimum height=1.3cm, align=center, font=\small, inner sep=6pt},
    sidebox/.style={rectangle, draw=black, rounded corners,
                fill=gray!15, minimum width=3.8cm,
                minimum height=0.85cm, align=center, font=\small},
    decision/.style={diamond, draw=black, fill=darkblue!15,
                     aspect=3, align=center, font=\small},
    arrow/.style={->, >=stealth, thick},
    label/.style={font=\small, midway},
]

\def\sep{1.55}

% ---- Nœuds principaux ----
\node[box] (dnin)   at (0,  0)
    {$\Delta n_\text{in}^{(i)}(\rv)$};
\node[font=\scriptsize\bfseries, anchor=north west]
    at (dnin.north west) {(1)};

\node[box] (dH)     at (0, -1.3*\sep)
    {$\Delta\hat{H}_\text{KS}
       = \Delta V_\text{PP}
       + \displaystyle\int d^3\rv'\,\Delta n(\rv')\,
         \frac{\delta\!\left[V_\text{H}+V_\text{XC}\right]\!(\rv)}
              {\delta n(\rv')}$};
\node[font=\scriptsize\bfseries, anchor=north west]
    at (dH.north west) {(3)};

\node[box] (stern)  at (0, -2.7*\sep)
    {Sternheimer projetée :\\[2pt]
     $\left(\hat{H}_\text{KS}-\varepsilon_{n\kv}\right)
      \ket{\Delta\psi_{n\kv}}
      = -\hat{P}_c\,\Delta\hat{H}_\text{KS}\ket{\psi_{n\kv}}$};
\node[font=\scriptsize\bfseries, anchor=north west]
    at (stern.north west) {(5)};

\node[box] (dnout)  at (0, -4.1*\sep)
    {$\Delta n_\text{out}(\rv)
       = 4\Re\displaystyle\sum_{n\kv} f_{n\kv}\,
         \psi^*_{n\kv}(\rv)\,\Delta\psi_{n\kv}(\rv)$};
\node[font=\scriptsize\bfseries, anchor=north west]
    at (dnout.north west) {(6)};

\node[box] (mix)    at (0, -5.3*\sep)
    {Mixing $\;\longrightarrow\;\Delta n_\text{in}^{(i+1)}(\rv)$};
\node[font=\scriptsize\bfseries, anchor=north west]
    at (mix.north west) {(7)};

\node[decision] (conv) at (0, -6.5*\sep) {Convergé?};

\node[box] (cel)    at (0, -8.0*\sep)
    {$\widetilde{C}^\text{el}_{\alpha\mu,\alpha'\mu'}(\qv)
      = \displaystyle\int d^3\rv\,\Delta n(\rv)\,\Delta' V_\text{PP}(\rv)
      + \displaystyle\int d^3\rv\,n(\rv)\,\Delta\Delta' V_\text{PP}(\rv)$};
\node[font=\scriptsize\bfseries, anchor=north west]
    at (cel.north west) {(8)};

\node[box] (dyn)    at (0, -9.4*\sep)
    {$D_{\alpha\mu,\alpha'\mu'}(\qv)
      = \dfrac{1}{\sqrt{M_\alpha M_{\alpha'}}}\,
        \widetilde{C}_{\alpha\mu,\alpha'\mu'}(\qv)$};
\node[font=\scriptsize\bfseries, anchor=north west]
    at (dyn.north west) {(9)};

\node[box] (freq)   at (0, -10.8*\sep)
    {Fréquences $\omega_{\nu\qv}$ et polarisations $\mathbf{e}^{(\nu)}(\qv)$};
\node[font=\scriptsize\bfseries, anchor=north west]
    at (freq.north west) {(10)};

% ---- Nœuds latéraux (précalculés) ----
\node[sidebox] (dVPP) at (6.8, -1.3*\sep)
    {$\Delta V_\text{PP}(\rv)$\\(précalculé)};
\node[font=\scriptsize\bfseries, anchor=north west]
    at (dVPP.north west) {(2)};

\node[sidebox] (gs)   at (6.8, -2.7*\sep)
    {$\psi_{n\kv},\;\varepsilon_{n\kv},\;n(\rv)$\\(calcul SCF préalable)};
\node[font=\scriptsize\bfseries, anchor=north west]
    at (gs.north west) {(4)};

% ---- Flèches flux principal ----
\draw[arrow] (dnin)  -- (dH);
\draw[arrow] (dH)    -- (stern);
\draw[arrow] (stern) -- (dnout);
\draw[arrow] (dnout) -- (mix);
\draw[arrow] (mix)   -- (conv);
\draw[arrow] (conv)  -- (cel)
    node[label, right] {Oui};
\draw[arrow] (cel)   -- (dyn)
    node[label, right] {$+\;\widetilde{C}^\text{ion}$};
\draw[arrow] (dyn)   -- (freq)
    node[label, right] {Diagonalisation};

% ---- Flèches latérales ----
\draw[arrow] (dVPP)  -- (dH);
\draw[arrow] (gs)    -- (stern);

% ---- Flèche de retour (boucle autocohérente) ----
\draw[arrow] (conv.west)
    -- ++(-2.9, 0)
    |- (dnin.west)
    node[pos=0.01, right, font=\small] {Non};

% ---- Rectangle délimitant la boucle ----
\draw[dashed, gray, rounded corners]
    (-4.6, 1.2) rectangle (4.6, -5.85*\sep);
\node[font=\small\itshape, gray] at (-3.0, 0.92)
    {Boucle DFPT};

% ---- Note : cycle répété pour chaque perturbation ----
\node[font=\scriptsize\itshape, gray, align=left, anchor=west]
    at (5.2, -5.3*\sep)
    {Boucle DFPT répétée pour \\chaque perturbation $(\qv\alpha\mu)$\\
     et chaque $\qv\in$ IBZ.\\[2pt]
     Notation : $\Delta\equiv\Delta_{\qv\alpha\mu}$,\\
     $\Delta'\equiv\Delta_{\qv\alpha'\mu'}$};

\end{tikzpicture}